\section*{Bab \ref{ch:b3:nervous-systems} --- Tatanan Sistem Saraf}

\begin{solution}{exo:b3:nervous-systems:1}
\omterm{def:b3:nervous-systems:cells}{Astrosit}: menyangga kalium dan glutamat, memberi makan neuron, memicu
penghadang darah--otak, mengatur aliran darah setempat. \omterm{def:b3:nervous-systems:cells}{Oligodendrosit}:
memielinkan akson di dalam sistem saraf pusat. Sel Schwann:
memielinkan akson tepi (dan memandu pemulihannya). \omterm{def:b3:nervous-systems:cells}{Mikroglia}:
\omterm{def:b3:innate-immunity:innate}{makrofag} otaknya, yang memangkas sinapsis dan membersihkan puing.
(Sel ependima melapisi bilik otaknya lalu membuat \omterm{prop:b3:nervous-systems:barrier-budget}{cairan serebrospinal}.)
\end{solution}

\begin{solution}{exo:b3:nervous-systems:2}
Pewarnanya menghitamkan beberapa neuron seluruhnya lalu membiarkan selebihnya
tak terlihat, sehingga sebuah sel tunggal dapat dilihat utuh. Cajal menyimpulkan bahwa
neuron adalah sel terpisah yang bersentuhan satu sama lain tanpa melebur, dengan
arah aliran yang tetap dari dendrit ke akson; Golgi menyimpulkan bahwa
mereka membentuk sebuah jaringan sinambung. Cajal benar; mikroskop
elektron menunjukkan celah sinapsisnya pada 1954.
\end{solution}

\begin{solution}{exo:b3:nervous-systems:3}
Sebuah ketukan meregangkan kuadrisepnya; kumparan ototnya menyala; aferen
Ia memasuki sumsumnya lalu bersinapsis langsung pada neuron motorik
kuadrisepnya, yang menyala lalu mengerutkan ototnya, sementara sebuah \omterm{def:b3:nervous-systems:cells}{interneuron}
menghambat neuron motorik otot paha belakangnya. Satu sinapsis, akson pendek pada
\qty{80}{m/s}: sekitar \qty{30}{ms}. Sebuah tendangan sukarela melewati
otaknya --- persepsi, keputusan, perencanaan gerak, berkas menurun,
beberapa sinapsis --- lalu memakan \qtyrange{200}{300}{ms}.
\end{solution}

\begin{solution}{exo:b3:nervous-systems:4}
Simpatis (noradrenalin pada sasarannya): jantung lebih cepat dan lebih kuat,
biji mata melebar, gerak dan sekresi ususnya berkurang, saluran napasnya melebar.
Parasimpatis (asetilkolin): jantung melambat, biji mata menyempit, gerak
dan sekresi ususnya bertambah, saluran napasnya menyempit.
\end{solution}

\begin{solution}{exo:b3:nervous-systems:5}
$\lambda = \sqrt{R_{m}d/4R_{i}} = \sqrt{10^{4}\times 10^{-4}/400}$ cm
$= \qty{0.05}{cm} = \qty{0.5}{mm}$. Sesudah \qty{2}{mm}: $e^{-4} \approx
2\,\%$. Melipatduakan $\lambda$ memerlukan empat kali garis tengahnya, yakni
\qty{4}{\micro m}.
\end{solution}

\begin{solution}{exo:b3:nervous-systems:6}
Indra: $12\times 6 = \qty{72}{m/s}$, $1.2/72 = \qty{17}{ms}$. Motorik:
\qty{90}{m/s}, $1.1/90 = \qty{12}{ms}$. Dua sinapsis, \qty{1}{ms}: sekitar
\qty{30}{ms}. Tak bermielin: $1.1\sqrt{12} = \qty{3.8}{m/s}$ dan
$1.1\sqrt{15} = \qty{4.3}{m/s}$: $315 + 258 + 1 \approx \qty{570}{ms}$
--- terlalu lambat untuk menyelamatkan kakinya.
\end{solution}

\begin{solution}{exo:b3:nervous-systems:7}
$2.4\times 10^{20}/8.6\times 10^{10} \approx 2.8\times 10^{9}$ ATP tiap
neuron tiap detik; pada $2\times 10^{9}$ tiap lecutan, sekitar \qty{1.4}{Hz}
rata-rata. Sandinya harus jarang: sedikit lecutan, dengan keterangan diusung oleh
neuron mana yang menyala dan kapan, bukan oleh laju tinggi di mana-mana.
\end{solution}

\begin{solution}{exo:b3:nervous-systems:8}
Dua neuron yang saling menghambat tanpa lelah membentuk sebuah sakelar: yang
pertama menyala menekan yang lain selama dorongannya bertahan, dan
tidak pernah ada serah terima. Pengayunannya memerlukan sebuah proses lambat yang melemahkan
cengkeraman pemenangnya --- penyesuaian penyalaannya, penekanan sinapsis
penghambatnya, atau lentingan yang kalah sesudah terlepas --- yang
tetapan waktunya menetapkan kalanya.
\end{solution}

\begin{solution}{exo:b3:nervous-systems:9}
L-dopa diusung melintasi penghadangnya oleh pengangkut bagi asam amino netral
yang besar; dopamin, yang polar dan bermuatan, tidak. Pasien meminum
L-dopa (dengan sebuah penghambat enzim tepinya sehingga ia tidak
terubah sebelum menyeberang), lalu enzim otaknya sendiri membuat dopamin
darinya di tempat ia diperlukan. \omterm{def:b3:adaptive-immunity:antibody}{Antibodi}, pada \qty{150}{kDa}, tidak menyeberangi
sambungan rapatnya maupun pengangkut mana pun; penghantarannya memerlukan sebuah
ulang-alik yang membajak sebuah reseptor untuk transitosis, atau suntikan ke dalam
\omterm{prop:b3:nervous-systems:barrier-budget}{cairan serebrospinalnya}.
\end{solution}

\begin{solution}{exo:b3:nervous-systems:10}
Arus selaput tiap satuan panjangnya adalah $V/r_{m}$ yang resistif ditambah
$c_{m}\,\partial V/\partial t$ yang kapasitif; kekekalan muatan memberi
$\partial I/\partial x = -(V/r_{m} + c_{m}\,\partial V/\partial t)$, dan
dengan $\partial V/\partial x = -r_{i}I$, $\frac{1}{r_{i}}\,\partial^{2}V/
\partial x^{2} = V/r_{m} + c_{m}\,\partial V/\partial t$; dengan mengalikannya dengan
$r_{m}$: $\lambda^{2}\,\partial^{2}V/\partial x^{2} = V + \tau\,\partial
V/\partial t$. Dengan $\partial V/\partial t = 0$ persamaan teoremanya
kembali. $\lambda$ sebuah panjang dan $\tau$ sebuah waktu, sehingga $\lambda/\tau$
sebuah laju --- yakni skala alami bagi seberapa cepat sebuah gangguan menyebar.
\end{solution}

\begin{solution}{exo:b3:nervous-systems:11}
$v \propto \sqrt{d}$: empat kali lajunya memerlukan enam belas kali
garis tengahnya, yakni \qty{8}{mm}. Luasnya $\pi(4\,\text{mm})^{2} \approx \qty{50}{mm^2}$
terhadap $\pi\,(10\,\unit{\micro m})^{2} = \qty{314}{\micro m^2}$: yakni nisbah
$1.6\times 10^{5}$. Sebuah saraf berisi serat tak bermielin yang cepat mustahil
--- satu serat semacam itu akan setebal sarafnya --- sehingga seekor hewan dengan
banyak saluran cepat memerlukan \omterm{def:b3:nervous-systems:cells}{mielin}, dan itulah sebabnya vertebrata (dan beberapa
hewan tak bertulang belakang secara mandiri) mengevolusikannya.
\end{solution}

\begin{solution}{exo:b3:nervous-systems:12}
$a = 0.4/30^{0.75} = 0.4/12.8 = 0.031$; bagi \qty{70000}{g}: $0.031\times
\num{70000}^{0.75} = 0.031\times 4300 \approx \qty{134}{g}$. Manusia yang
\qty{1350}{g} adalah sepuluh kali ramalannya. Massa otak mengikuti massa
tubuh karena tubuh yang lebih besar memerlukan lebih banyak neuron indra dan motorik, dan
karena neuronnya sendiri tumbuh seiring ukuran otaknya; yang menjejaki kognisinya
adalah cacah neuron di dalam korteksnya, yang bergantung pada seberapa padat
mereka berkemas --- primata berkemas lebih banyak tiap gram daripada mamalia lain,
dan manusia paling banyak dari semuanya.
\end{solution}

\begin{solution}{pb:b3:nervous-systems:1}
\textbf{1.} $\lambda = \sqrt{2\times 10^{4}\,d/400}$: bagi $d =
\qty{1e-4}{cm}$, $\sqrt{5\times 10^{-3}} = \qty{0.071}{cm} =
\qty{0.71}{mm}$; bagi \qty{4}{\micro m}, \qty{1.41}{mm}. $\tau = 2\times
10^{4}\times 10^{-6} = \qty{20}{ms}$.
\textbf{2.} $5\,e^{-0.3/0.71} = \qty{3.3}{mV}$; $5\,e^{-0.3/1.41} =
\qty{4.0}{mV}$.
\textbf{3.} Masukan yang jauh tiba dengan teredam, sehingga tempat sebuah sinapsis mendarat
menetapkan bobotnya; sebuah potensial meluruh dalam sekitar $\tau$, sehingga dua masukan
terjumlahkan hanya bila keduanya jatuh dalam sekitar \qty{20}{ms} satu sama lain --- neuronnya
adalah sebuah pendeteksi kebetulan dengan jendela \qty{20}{ms}.
\textbf{4.} Bermielin \qty{60}{m/s}; tak bermielin $1.1\sqrt{10} =
\qty{3.5}{m/s}$; untuk mencapai \qty{60}{m/s} tanpa \omterm{def:b3:nervous-systems:cells}{mielin}, $d = (60/1.1)^{2}
\approx \qty{3000}{\micro m}$: yakni tiga milimeter.
\textbf{5.} $\pi(5)^{2} = \qty{79}{\micro m^2}$ terhadap $\pi(1500)^{2} =
7\times 10^{6}\,\unit{\micro m^2}$: yakni sekitar \num{90000} akson bermielin
di dalam ruang satu akson itu.
\textbf{6.} \qty{2}{cm} pada \qty{60}{m/s} adalah \qty{0.3}{ms}; pada
\qty{3.5}{m/s}, \qty{5.7}{ms}: yakni sekitar \qty{5}{ms} penundaan tambahan. Lebih buruk lagi,
selaput telanjangnya punya kapasitans besar dan sedikit saluran, sehingga
arus dari simpul utuh yang terakhir dapat gagal membawanya ke ambangnya:
itulah blok hantaran.
\textbf{7.} Indra $1.2/60 = \qty{20}{ms}$; dua sinapsis \qty{1}{ms};
motorik $1.1/90 = \qty{12}{ms}$: yakni sekitar \qty{33}{ms} (ditambah semilidetik
pada sambungan saraf--ototnya).
\textbf{8.} Ke otaknya: $0.6/60 = \qty{10}{ms}$ ditambah dua sinapsis,
yakni \qty{11}{ms} sesudah isyaratnya memasuki sumsumnya pada \qty{20}{ms}:
isyaratnya mencapai korteksnya kira-kira ketika kakinya bergerak;
pengalaman sadar nyerinya memerlukan beberapa ratus milidetik lagi
pengolahan, sehingga kakinya ditarik sebelum nyerinya terasa.
\textbf{9.} $1.1\sqrt{1} = \qty{1.1}{m/s}$: $1.8/1.1 \approx
\qty{1.6}{s}$. Nyeri pertama yang tajam tiba lewat serat bermielin tipis
dalam puluhan milidetik; nyeri kedua yang tumpul dan membakar lewat
serat C tak bermielin sedetik atau lebih kemudian.
\textbf{10.} Indra $3.2/60 = \qty{53}{ms}$, motorik $3.1/90 =
\qty{34}{ms}$, sinapsisnya \qty{1}{ms}: sekitar \qty{90}{ms}. Hewan yang panjang
punya refleks yang lambat; mereka mengimbanginya dengan akson yang lebih tebal dan dengan
mewakilkan lebih banyak kepada sirkuit setempat.
\textbf{11.} $0.8/80 = \qty{10}{ms}$ tiap arah, yakni \qty{20}{ms}; sinapsisnya
\qty{0.5}{ms}, sambungannya \qty{1}{ms}, gayanya \qty{5}{ms}: yakni \qty{26.5}{ms},
dengan transduksi kumparannya sendiri dan hantaran ke dalam ototnya
melengkapi sisa \qty{30}{ms} itu.
\textbf{12.} Akson tak bermielin menghantar pada semeter atau dua tiap detik,
sehingga gelung yang memakan \qty{30}{ms} pada orang dewasa memakan ratusan, dengan
ketepatan waktu yang buruk; pemielinannya menaikkan lajunya sepuluh sampai lima puluh kali lipat lalu
membuat ketepatan waktunya cermat, dan itulah yang dituntut berjalan, menggenggam, dan berbicara.
\textbf{13.} $20/50\,000 = 4\times 10^{-4}$ mol/s $= 2.4\times 10^{20}$
ATP tiap detik; $2.8\times 10^{9}$ tiap neuron tiap detik.
\textbf{14.} $4\times 10^{8} + 7000\times 2\times 10^{4} = 4\times 10^{8}
+ 1.4\times 10^{8} = 5.4\times 10^{8}$ ATP.
\textbf{15.} Semuanya untuk lecutan: $2.8\times 10^{9}/5.4\times 10^{8} \approx
\qty{5}{Hz}$; separuhnya: sekitar \qty{2.6}{Hz}.
\textbf{16.} Pada \qty{1}{Hz}, $5.4\times 10^{8}/2.8\times 10^{9} \approx
\qty{20}{\%}$ anggarannya: sandinya jarang --- kebanyakan neuronnya bungkam
hampir sepanjang waktu, dengan keterangan pada neuron mana yang sedikit menyala dan kapan.
\textbf{17.} $8.6\times 10^{10}\times 50\times 5.4\times 10^{8} = 2.3
\times 10^{21}$ ATP/s $= 3.9\times 10^{-3}$ mol/s $\approx \qty{190}{W}$
--- yakni dua kali seluruh tubuh yang beristirahat.
\textbf{18.} $5\times 16 = \qty{80}{kJ/h} = \qty{22}{W}$: pas-pasan cukup,
tanpa cadangan. Ketika pasokannya berhenti, pompanya berhenti dalam hitungan detik,
landaiannya habis, dan kesadarannya hilang dalam sekitar sepuluh detik.
\textbf{19.} $a = 0.4/12.8 = 0.031$; \qty{70}{kg}: $0.031\times 4300 =
\qty{134}{g}$; \qty{5000}{kg}: $0.031\times 1.06\times 10^{5} \approx
\qty{3300}{g}$.
\textbf{20.} Manusia $1350/134 \approx 10$; gajah $4800/3300 \approx
1.5$.
\textbf{21.} Di dalam \omterm{def:b3:nervous-systems:plan}{otak kecilnya} --- \qty{98}{\%} di antaranya --- yang menyelaraskan
empat puluh ribu otot belalainya. Cacah neuron korteksnya, bukan totalnya,
yang menjejaki kognisinya; \omterm{def:b3:nervous-systems:plan}{otak kecilnya} menjejaki tugas motorik menjalankan
sebuah tubuh raksasa.
\textbf{22.} Akson dan dendritnya memanjang seiring otaknya, sehingga tiap neuron
menempati volume lebih besar dan cacah neuronnya tumbuh lebih lambat daripada
massanya. Pada primata ukuran neuronnya nyaris tetap sementara otaknya tumbuh,
sehingga cacah neuronnya tumbuh hampir sebanding dengan massanya, dan sebuah otak
primata pada bobot tertentu memuat beberapa kali neuron sebuah otak hewan
pengerat pada bobot itu.
\textbf{23.} Sebuah lengan dapat menyelaraskan pengisap dan sendinya sendiri, menjangkau,
menggenggam, mengoper makanan sepanjang dirinya, dan menghindari menggenggam kulitnya sendiri, lewat
sirkuit setempat; ia tidak dapat melihat, memilih sebuah sasaran, atau mempelajari sebuah pembedaan.
Ujinya: potonglah tali saraf lengannya dari otaknya lalu sentuhlah lengannya dengan
makanan --- ia menggenggam lalu mengopernya ke arah tempat mulutnya berada; sebuah
lengan yang teramputasi melakukan hal yang sama selama sejam.
\textbf{24.} Sekitar $23$ sinapsis tiap neuron terhadap $7000$: yakni tiga ratus
kali lebih sedikit. Petanya yang lengkap mengatakan neuron mana yang dapat memengaruhi neuron mana,
tetapi bukan kekuatan atau tanda tiap sambungannya, neuromodulator mana yang
mengubahnya, atau dinamikanya --- sehingga bahkan bagi cacingnya, diagram
perangkaiannya meramalkan tingkah lakunya hanya sebagian.
\textbf{25.} $\lambda = \qty{0.71}{mm}$ bagi dendrit \qty{1}{\micro m};
selang waktu penarikannya sekitar \qty{33}{ms}; laju penyalaan rata-ratanya
sekitar \qty{2.6}{Hz} dengan separuh anggarannya untuk lecutan.
\end{solution}
